% Methods — see paper/methods_skeleton.md for the long-form outline.
% Target 1500-2000 words for PLOS Comp Bio. Each subsection below maps to
% a markdown subsection in methods_skeleton.md.

\subsection{Model overview}\label{sec:methods:overview}
The two-leg spinal CPG presented here extends a canonical computational
sCPG model~\cite{rybak2006,shevtsova2015} with four
mechanisms not present in the Zhang 2022 reference model~\cite{zhang2022}:
(i) STDP-driven self-organisation of descending and cutaneous afferent
synapses onto the rhythm generator;
(ii) a closed-loop Ia feedback pathway to the reciprocal-inhibition
interneurons;
(iii) motoneuron pools driving Hill-like muscle proxies with force and
length dynamics;
(iv) closed-loop, force-triggered stance/swing switching: each leg's
cutaneous (paw-contact) input is switched on and off by that leg's own
extensor force rather than by a clock, with sequential Ia activation
simulating heel-to-toe weight transfer within each stance.
The model is implemented in NEST~3.9~\cite{nest2002} and is publicly
available at \href{https://github.com/max-talanov/tinyCPG}{github.com/max-talanov/tinyCPG}.

\begin{figure}[t]
  \centering
  \includegraphics[width=0.8\textwidth]{fig1_schematic.png}
  \caption{\textbf{Architecture of the two-leg spinal CPG.}
    Each leg contains a flexor (F)
    and extensor (E) rhythm-generator half-centre coupled by asymmetric
    reciprocal inhibition via InE and InF interneurons (Zhang 2022 6:1
    ratio). Sensory afferents (Ia, cutaneous) feed into the
    reciprocal-inhibition interneurons; motor pools (M-E, M-F) drive
    muscle proxies whose force and length signals close the Ia loop.
    L\,$\leftrightarrow$\,R commissural inhibition (V0-like) is mediated by
    flexor- and extensor-targeted projections.}
  \label{fig:schematic}
\end{figure}

\subsection{Single-neuron model}\label{sec:methods:neuron}
The network is implemented at two complementary levels of biophysical
detail, using the same architecture, the same plastic projections and the
same experimental protocol in both cases.

\paragraph{Reduced (Izhikevich) implementation.}
In the NEST implementation, all spiking neurons are modelled as Izhikevich
quadratic-integrate-fire units~\cite{izhikevich2003}. Each neuron obeys
\begin{align}
  \frac{dv}{dt} &= 0.04v^2 + 5v + 140 - u + I, \\
  \frac{du}{dt} &= a\,(bv - u),
\end{align}
with the after-spike reset $v \leftarrow c$, $u \leftarrow u + d$ applied
whenever $v \geq 30$\,mV. Excitatory extensor RG-E neurons use the
regular-spiking set ($a=0.02, b=0.2, c=-65, d=8$). Flexor RG-F neurons use
the intrinsically-bursting set ($a=0.02, b=0.2, c=-55, d=4$), which
approximates persistent-sodium-driven ($I_\text{NaP}$) bursting
phenomenologically rather than deriving it from channel kinetics.
Inhibitory interneurons (InE, InF, Ia-INT) use a fast-spiking set
($d=2$). This formulation reproduces the spiking and bursting regimes
required by the rhythm generator at a small fraction of the computational
cost of a conductance-based model, which is what makes the large
mode\,$\times$\,learning-rate parameter sweeps reported here tractable.

\paragraph{Conductance-based (Hodgkin--Huxley) implementation.}
Because the reduced model represents bursting phenomenologically, we also
implemented the same circuit with conductance-based neurons in
NEURON~\cite{hines1997}, following the Hodgkin--Huxley
formalism~\cite{hodgkinHuxley1952}. Membrane potential evolves as
\begin{equation}
  C_m \frac{dV}{dt} = -\!\sum_i g_i\, m_i^{p}h_i^{q}\,(V - E_i) - I_\text{leak} + I_\text{syn},
  \label{eq:hh}
\end{equation}
where each active conductance $g_i$ is gated by activation ($m_i$) and
inactivation ($h_i$) variables obeying first-order kinetics
$\dot{x} = (x_\infty(V) - x)/\tau_x(V)$. Rhythm-generator neurons include
fast sodium, delayed-rectifier potassium and leak conductances together
with an explicit persistent sodium current ($I_\text{NaP}$), so that
bursting arises from channel kinetics rather than from a reset rule.
Motoneurons and interneurons are multi-compartment cells with somatic and
dendritic sections. This implementation resolves the sub-threshold and
channel-level dynamics that the Izhikevich formulation abstracts away, at
substantially higher computational cost; it therefore serves as a
biophysical cross-check on the reduced model rather than as the vehicle
for the full parameter sweep.%
\todo{The NEURON/HH results are not yet included in Results. Until they
are, this subsection describes an implementation the paper does not yet
report on, and the Limitations no longer flag the absence of a
conductance-based model --- add the HH results (or re-scope this
paragraph) before circulating.}

Synaptic delays in both implementations are drawn from a rat
\texttt{length\_velocity} preset (Table~\ref{tab:delays}) with
$\sigma=0.2$\,ms Gaussian jitter.

\subsection{Network architecture}\label{sec:methods:network}

\paragraph{Asymmetric reciprocal inhibition (Zhang 2022).}
The F$\to$InF$\to$E pathway is strong ($W_{\text{InF}\to\text{RG-E}}=-48$\,pA,
$P=0.30$); the E$\to$InE$\to$F pathway is weak
($W_{\text{InE}\to\text{RG-F}}=-8$\,pA, $P=0.15$).
The 6:1 ratio preserves the flexor's role as the rhythm-leading
half-centre~\cite{zhang2022,talpalar2013}.

\paragraph{Closed-loop Ia$\to$reciprocal-IN pathway (this work).}
Ia-E parrots project to InE ($W=6$\,pA, $P=0.25$) and Ia-F parrots to
InF, providing closed-loop sensory drive into the reciprocal-inhibition
core. This pathway compensates the loss of strong tonic descending drive
characteristic of deafferented fictive locomotion~\cite{pearson1995,
rossignol2006,hultborn2006}.

\paragraph{Cutaneous reflex.}
Cutaneous parrots project to RG-E (STDP-plastic) and to InE
($W=6$\,pA, $P=0.30$), implementing the stance-phase cutaneous-reflex
gating of extensor activity~\cite{forssberg1980,pearson1995}.

\paragraph{Commissural inhibition.}
L$\leftrightarrow$R flexor RG inhibition is strong
($W=-20$\,pA, $P=0.22$); extensor inhibition is weak
($W=-8$\,pA, $P=0.10$). This V0-mediated commissural inhibition
enforces left--right alternation~\cite{kiehn2016,talpalar2013}.

\paragraph{Motor pool, muscle relay.}
Each motor neuron projects to $N_{\text{mus}}=100$ muscle-parrot relays
($P_{\text{M}\to\text{mus}}=0.8$). Motor pool reciprocal inhibition
($W=-22$\,pA, $P=0.25$) enforces antagonist suppression at the motor
level~\cite{mccreaRybak2008}.

\paragraph{Connectivity specification.}
The full per-connection weight and delay distributions of every
projection were obtained by querying NEST directly after network
construction. Table~\ref{tab:delays}
lists all projections grouped by the functional blocks of the
architecture (Fig.~\ref{fig:schematic}): descending/supraspinal drive,
the sensory afferent$\to$RG projections, the rhythm-generator reciprocal
core, the Ia proprioceptive interneuron pathway, motor output, and the
commissural (interlimb) projections. Static projections are
heterogeneous: their per-connection weights are drawn from a lognormal
distribution about the design value (coefficient of variation $0.5$, mean
and sign preserved), since biological synaptic weights are lognormally
distributed rather than uniform~\cite{song2005,buzsaki2014}. The mean
preserves the designed balance---in particular the asymmetric reciprocal
inhibition (InF$\to$RG-E $=-48$\,pA, InE$\to$RG-F $=-8$\,pA) reproduces the
Zhang~6:1 ratio exactly. The plastic projections (CUT$\to$RG-E and
BS$\to$RG-\{E,F\}) start from a lognormal distribution and are reported
both as initialised and after learning (Fig.~\ref{fig:connectivity}a,b;
Table~\ref{tab:delays}). The cutaneous input to RG-E is a single plastic projection. Conduction delays follow the rat
\texttt{length\_velocity} preset (brainstem pathways $\approx 3.0$\,ms,
reciprocal and motor pathways $\approx 1.8$\,ms, commissural
$\approx 2.2$\,ms) with $0.2$\,ms across-connection jitter. The
muscle$\to$Ia spindle/GTO feedback is implemented as a rate-coded
transduction (force and length $\to$ Ia firing rate), not a synaptic
projection, and so has no weight/delay entry.
Figure~\ref{fig:connectivity} shows the weight and delay distributions.


\begin{figure}[t]
  \centering
  \includegraphics[width=0.8\textwidth]{fig_connectivity.png}
  \caption{\textbf{Synaptic weights and conduction delays along every
  projection, before and after learning.} \textbf{(a)} Weight histograms of
  the three plastic projections, leg L: as initialised (grey, lognormal
  $\mu=3.5$\,pA, CV$=0.3$) and after 120\,s of learning (orange; medium walk,
  $\lambda=10^{-3}$). \textbf{(b)} Synaptic weight per projection (mean $\pm$
  per-connection s.d.; green excitatory, red inhibitory); plastic
  projections are shown at their learned values (open circle = initial
  mean). All static projections are heterogeneous (lognormal, CV$\approx
  0.5$), with the asymmetric reciprocal inhibition InF$\to$RG-E $=-48$\,pA
  versus InE$\to$RG-F $=-8$\,pA preserving the Zhang~6:1 ratio in the mean.
  \textbf{(c)} Conduction~+~synaptic delay per projection (mean $\pm$
  across-connection s.d.), following the rat \texttt{length\_velocity}
  preset with $0.2$\,ms jitter.}
  \label{fig:connectivity}
\end{figure}

\subsection{Brainstem drive}\label{sec:methods:bs}
Reticulospinal input is modelled as a tonic Poisson process at
$BS_{\text{regular}}=60$\,Hz with $\sigma=0.25$\,Hz noise, identical for
both legs. The 60\,Hz value sits in the
middle of the rat reticulospinal range
(20--80\,Hz)~\cite{drew1986,brocard2010}. In contrast to~\cite{zhang2022},
who vary tonic drive $\alpha$ to vary gait speed, we hold the BS firing
rate tonic and let the stance/swing timing emerge from sensory feedback
(see \autoref{sec:methods:paced}); the \emph{strength} of the BS$\to$RG
synapses is learned (\autoref{sec:methods:stdp}).

\subsection{Spike-timing-dependent plasticity}\label{sec:methods:stdp}
Two projections --- cutaneous CUT$\to$RG-E and brainstem
BS$\to$RG-\{E,F\} --- carry multiplicative STDP synapses (NEST
\verb|stdp_synapse|; $\tau_+=20$\,ms, default $\lambda=10^{-3}$,
$\alpha=0.95$, $\mu_+=\mu_-=0.4$, $w_{\max}=120$\,pA for CUT and
$w_{\max}=30$\,pA for BS)~\cite{biPoo1998,morrison2007}. Initial
weights are drawn from a lognormal distribution with mean $\mu$ and
coefficient of variation CV ($\mu=3.5$\,pA, CV$=0.30$ unless stated
otherwise). We swept $(\mu,\text{CV})$ across ten diagnostic points
($\mu$ from $0$ to $16$\,pA) to verify that the learned weights converge
to the same plateau regardless of initialisation (see
Section~\ref{sec:results:robustness}). The proprioceptive
Ia$\to$RG projections are not plastic in this configuration; Ia
feedback acts through the static reciprocal-interneuron and
Ia-interneuron pathways only.

\paragraph{Extensor-biased plasticity and stance/swing asymmetry.}
The plastic and paced drives are deliberately extensor-weighted, mirroring
the stance/swing organisation of locomotor control. The cutaneous
afferent --- the only projection with the large ($w_{\max}=120$\,pA)
plasticity ceiling --- targets the extensor (CUT$\to$RG-E) and has no
flexor counterpart, reflecting that paw-contact cutaneous feedback fires
during weight-bearing stance and reinforces the load-bearing
extensor~\cite{pearson1995,rossignol2006}. Likewise the sequential
heel$\to$toe Ia-E drive reinforces the extensor within each stance, whereas the
flexor receives no external pacing and is generated intrinsically ---
by the persistent-sodium-like bursting of the RG-F units and by release
from reciprocal inhibition during swing. As a consequence the extensor
half-centre is sensory-driven and tightly clocked while the flexor is
centrally generated and intrinsically more variable; and because the
dominant plastic projection is extensor-only, increasing $\lambda$
strengthens the descending drive asymmetrically, in favour of the
extensor (analysed in Section~\ref{sec:results:force}).

\paragraph{Learning-rate variation.}
For the locomotion-mode $\times$ learning-rate grid
(\autoref{sec:methods:protocols}), the STDP learning rate $\lambda$ was
varied across five decades around the cortical bio-plausible STDP
rate~\cite{biPoo1998,morrison2007}:
$\lambda \in \{10^{-2},10^{-3},10^{-4},10^{-5},10^{-6}\}$. With a
120\,s run the two slowest rates do not converge, which lets us probe the
regime of weak, still-developing learned weights. The parameter is
recorded to the simulation output for every run.

\subsection{Sensory feedback and Ia rate coding}\label{sec:methods:ia}
Ia afferent firing rate $r_{\text{Ia}}$ is updated every 100\,ms as a
linear combination of muscle force $F$ and stretch $\Delta\ell$:
\begin{equation}
  r_{\text{Ia}} = r_0 + k_F \cdot F + k_S \cdot \max(0, \ell - L_0),
  \quad \text{clamped to } [0, r_{\max}],
  \label{eq:ia}
\end{equation}
with $r_0=10$\,Hz, $k_F=6$\,Hz/F, $k_S=250$\,Hz/L,
$r_{\max}=500$\,Hz~\cite{loeb1981,prochazka1999}. This rate drives the
classical Ia-INT$\to$antagonist motor-pool
pathway~\cite{jankowska1992} and the closed-loop Ia$\to$recip-IN pathway
introduced in this work.

\paragraph{Graded sensory ablation (Courtine/Lavrov paradigm).}
To emulate the partial-weight-bearing and limb-unloaded preparations
used in rat SCI rehabilitation
experiments~\cite{lavrov2008,courtine2009,edgerton2008}, we expose a
multiplicative Ia-feedback gain $G_{\text{Ia}} \in [0, 1]$, applied
after Eq.~\ref{eq:ia}:
$\tilde{r}_{\text{Ia}} = G_{\text{Ia}}\cdot r_{\text{Ia}}$, and a matching
cutaneous gain $G_{\text{CUT}}$ that scales the stance-phase paw-contact
rate. Full weight-bearing uses $G_{\text{Ia}}=G_{\text{CUT}}=1.0$;
\emph{toe stepping} (partial weight-bearing) uses
$G_{\text{Ia}}=G_{\text{CUT}}=0.5$. The classical Ia-INT$\to$antagonist
pathway inherits the same gain, so the manipulation reduces \emph{all} Ia
output uniformly; the sequential heel$\to$toe Ia-E drive is not scaled.
Full unloading (\emph{air stepping}, $G=0.1$) is not reported: in pilot
simulations the extensor force at that loading was too small for the
force-triggered stance detector (\autoref{sec:methods:paced}) to separate
stance from fluctuations, and stance collapsed to single update steps, so
no closed-loop gait exists to analyse.

\subsection{Activation and force dynamics}\label{sec:methods:muscle}
The muscle relay rate $r_{\text{mus}}$ is mapped to activation $a$
through a saturating nonlinearity gated by RG-rate. The gate is a smooth
logistic (sigmoid) of the normalised RG population rate --- a
bio-plausible soft saturation in place of a hard clip:
\begin{align}
  a_{\text{raw}} &= a_{\max} \, \bigl(1 - \exp(-k_a \, r_{\text{mus}})\bigr), \label{eq:araw}\\
  d &= \sigma\!\Bigl(k_g\bigl(r_{\text{RG}}/r_{\text{ref}} - x_0\bigr)\Bigr),
       \qquad \sigma(z)=\frac{1}{1+e^{-z}}, \label{eq:gate}\\
  a_{\text{tgt}} &= a_{\text{raw}} \cdot d, \label{eq:atgt}\\
  a(t+\Delta t) &= a(t) + \bigl(1 - e^{-\Delta t/\tau_a}\bigr)\bigl(a_{\text{tgt}} - a(t)\bigr).
  \label{eq:alp}
\end{align}
with gate steepness $k_g=8$ and mid-point $x_0=0.6$. The logistic
saturates smoothly at $1$ and suppresses the inter-burst trough
($d\approx0.06$ at $r_{\text{RG}}/r_{\text{ref}}\approx0.25$ versus
$d\approx0.83$ at a burst peak), so the gated target $a_{\text{tgt}}$ is
bounded in $[0,a_{\max})$ by construction --- without the hard clip used
in earlier versions of the model.
Force follows an identical first-order low-pass driven by a saturating
function of activation:
\begin{align}
  F_{\text{tgt}} &= F_{\max} \, \bigl(1 - \exp(-k_F \, a)\bigr),\\
  F(t+\Delta t) &= F(t) + \bigl(1 - e^{-\Delta t/\tau_F}\bigr)\bigl(F_{\text{tgt}} - F(t)\bigr).
\end{align}
Under the gait protocol (\autoref{sec:methods:paced}) the time
constants are set to $\tau_a=40$\,ms, $\tau_F=80$\,ms. The latter
sits within the rat fast-twitch fibre rise-time
range~\cite{zajac1989,winters1995}.

\paragraph{Muscle fatigue.}
Each muscle carries a slow activity-dependent fatigue state
$\phi\in[0,\phi_{\max}]$ that scales its force ceiling,
$F_{\text{tgt}} = F_{\max}(1-\phi)\bigl(1-\exp(-k_F a)\bigr)$, and evolves as
\begin{equation}
  \frac{d\phi}{dt} = \frac{a}{a_{\max}}\,\frac{\phi_{\max}-\phi}{\tau_{\text{on}}}
    - \Bigl(1-\frac{a}{a_{\max}}\Bigr)\frac{\phi}{\tau_{\text{rec}}},
  \label{eq:fatigue}
\end{equation}
building up while the muscle is active and clearing while it is silent
($\phi_{\max}=0.95$; $\tau_{\text{on}}$ and $\tau_{\text{rec}}$ per
locomotion mode, Table~\ref{tab:modes}). Fatigue is what lets extensor force
decline during a sustained stance so that the force-triggered stance
detector can end it: the extensor rhythm generator has no intrinsic
burst-termination mechanism (only RG-F is intrinsically bursting), so
without fatigue the CUT$\to$RG-E$\to$force$\to$CUT loop holds force at its
ceiling indefinitely. It is the closest local analogue of
persistent-sodium-driven burst termination in this reduced model.

\paragraph{Force units and SI calibration.}
Force is reported in arbitrary units (a.u.). The dimensionless ceiling
$F_{\max}=25$\,a.u.\ in Eq.~(6) was chosen so the activation--force
cascade operates in a numerically convenient range; absolute Newton
calibration is not intrinsic to the present model and is conventional
in CPG-modelling reports~\cite{rybak2006,danner2017,zhang2022}. A
first-order linear rescaling maps a.u.\ to Newtons by anchoring on
\emph{in vivo} rat hindlimb data: the peak isometric tetanic force of
rat medial gastrocnemius is approximately
$5$--$7$\,N~\cite{roy1991,burkholderLieber2001}. Treating our reported
peak Force-E of $17.4$\,a.u.\ as representing peak ankle-extensor
tetanic force gives a conversion of $1$\,a.u.\ $\approx$
$0.30$--$0.40$\,N. Nothing in the model dynamics depends on this
choice of units; we use a.u.\ throughout for transparency.

\subsection{Force-triggered stance/swing protocol}\label{sec:methods:paced}
Stance and swing are not imposed by a clock. Each leg carries a stance
detector --- a Schmitt trigger on its own extensor force $F_E$ ---
evaluated every 100\,ms. During stance the detector tracks the running peak
$\hat F$ of the current bout; stance ends (``foot lifts'') when $F_E$ falls
below $\theta_{\text{off}}\hat F$, and the next stance begins (``foot
touches down'') when $F_E$ rises above $\theta_{\text{on}}\hat F$
($\theta_{\text{on}}=0.80$; $\theta_{\text{off}}$ per mode,
Table~\ref{tab:modes}). $\hat F$ is reset at each stance onset, so the
thresholds are relative to the leg's own current force scale. During
stance the cutaneous generator of that leg fires at
$G_{\text{CUT}}\cdot CUT_{\text{on}}$ ($CUT_{\text{on}}=100$\,Hz) and three
sequential Ia-E sub-groups fire at 60, 80 and 100\,Hz over consecutive
thirds of a nominal stance window, emulating heel$\to$mid$\to$toe pressure
transfer~\cite{loebDuysens2018}; during swing all CUT and Ia-E drive is
silenced and an 80\,Hz flexor afferent (hip/flexor stretch) drives RG-F.
Stance termination is therefore produced by the closed loop
CUT$\to$RG-E$\to$force, together with muscle fatigue
(Eq.~\ref{eq:fatigue}), which lets force decline within a bout.

Two timeouts bound each phase ($T_{\max}$, Table~\ref{tab:modes}), as an
endogenous-timer backstop analogous to the hip-extension limit that
triggers swing even under continued loading~\cite{grillnerRossignol1978}.
For stance this is a safeguard only: we report the fraction of stance
bouts that end at the timeout (``at cap'') and treat a condition as a
genuine closed loop only when this fraction is near zero. Swing, in
contrast, has no sensory termination signal in the model (there is no
hip-position feedback), so swing always lasts $T_{\max}$; cycle-to-cycle
variability and speed differences therefore arise from stance.

Left--right symmetry is broken at $t=0$ by starting the right leg in
stance; both legs' cutaneous input is on for an initial priming window
$T_{\text{lead}}$ so that both CUT$\to$RG-E synapses receive their first
potentiation. At the fast and toe-stepping operating points the two legs
otherwise settle into a phase relation that depends on spike-timing noise
(sometimes in-phase); there, a small persistent asymmetry in fatigue onset
($\tau_{\text{on}}(1\mp\epsilon)$ for the leading/lagging leg) keeps them
in alternation, reflecting that real limbs are not identical.

\begin{table}[t]
  \centering
  \footnotesize
  \caption{Force-triggered operating point per locomotion mode. $P$ sets
    the nominal stance window of the heel$\to$toe Ia-E sequence only; the
    actual stride is emergent. $G$: Ia and cutaneous feedback gain.
    Common to all modes: $\theta_{\text{on}}=0.80$, $\phi_{\max}=0.95$,
    100\,ms detector update. Slow walk is medium walk with all
    fatigue/timeout constants scaled by $1.3$. Values were screened at the
    production network size for a genuine closed loop (no stance bouts at
    the timeout) and L/R alternation.}
  \label{tab:modes}
  \begin{tabular}{l r r r r r r r r}
    \toprule
    Mode & $G$ & $P$ (ms) & $\tau_{\text{on}}$ (ms) & $\tau_{\text{rec}}$ (ms)
      & $\theta_{\text{off}}$ & $T_{\max}$ (ms) & $T_{\text{lead}}$ (ms) & $\epsilon$ \\
    \midrule
    slow walk     & 1.0 & 1200 & 340 & 780 & 0.35 & 585 & 150 & 0    \\
    medium walk   & 1.0 &  520 & 260 & 600 & 0.35 & 450 & 150 & 0    \\
    fast walk     & 1.0 &  600 & 156 & 360 & 0.30 & 380 &  90 & 0.04 \\
    toe stepping  & 0.5 &  730 & 260 & 438 & 0.35 & 450 & 110 & 0.12 \\
    \bottomrule
  \end{tabular}
\end{table}

\subsection{Simulation setup}\label{sec:methods:sim}
All simulations were performed with NEST~3.9~\cite{nest2002} using a
kernel resolution of 0.2\,ms and simulation chunks of 100\,ms, at the full
network size (Table~\ref{tab:delays}) with tonic BS at 60\,Hz. Grid runs
used 16 OpenMP threads per task on the MN5 general-purpose partition
(Barcelona Supercomputing Center) or 2--4 threads on a workstation; a
120\,s simulation completes in $\approx$10--25\,min wall-clock.
Multi-threaded NEST runs are not bit-reproducible across thread counts
(spike-delivery order varies), so results are reported as single runs per
condition with the initialisation sweep (\autoref{sec:results:robustness})
standing in for replication.

\subsection{Experimental protocols}\label{sec:methods:protocols}
Two in-silico experiments are reported. Each is encoded as a stand-alone
shell script that submits one SLURM job array (or runs locally) and writes
one HDF5 file per task. Pseudocode is given in the Supplementary material
(Algorithms~\ref{alg:main}--\ref{alg:grid}).

\paragraph{Initialisation robustness.}
Ten 120\,s runs of medium walk at $\lambda=10^{-3}$, one per STDP
initial-weight distribution $(\mu,\mathrm{CV})$ from the ten
initialisation points ($\mu=0$--$16$\,pA), seed $12345+10007\,i$
(\texttt{run\_cutforce\_sweep6.sh}). Used in
Section~\ref{sec:results:robustness}.

\paragraph{Locomotion-mode $\times$ learning-rate grid.}
Twenty 120\,s runs crossing the four locomotion modes of
Table~\ref{tab:modes} with the five learning rates, all at
$(\mu,\mathrm{CV})=(3.5,0.30)$, seed $12345$
(\texttt{run\_forcetrig\_grid.sh}). Used in
Sections~\ref{sec:results:weights}--\ref{sec:results:comparison}.

\subsection{Data analysis}\label{sec:methods:analysis}
Stance and swing bouts were taken from the logged per-step stance-detector
state. A condition's stance timing is \emph{genuine} (closed-loop) when the
fraction of stance bouts that end at the timeout $T_{\max}$ (``frac.\
stance at cap'', worst leg) is near zero; a high fraction means the timeout,
not the force threshold, is setting the rhythm. Counter-phase strength was
reported as Pearson's correlation between Force-E and Force-F of the same
leg, and left--right alternation as the correlation between the two legs'
Force-E. Peak force was reported as the 95th percentile to avoid outlier
sensitivity. Unless stated otherwise, metrics are computed over the last
20\,s of the run (Figs.~\ref{fig:mode_lambda_heat}--\ref{fig:mode_lambda_trend})
or from $t\ge30$\,s (initialisation sweep), after the learned weights have
converged. STDP convergence was reported as the per-projection mean
$\pm$ standard deviation of synaptic weights sampled every 1\,s. No formal
statistical tests are reported.

% --- Parameter tables go in supplementary ---
